% ============================================================
% Come si usano, argomento 8: variabili aleatorie multiple
% ============================================================
\argomento{Variabili aleatorie multiple}{Formule a p.~\pageref*{f:congiunta}}

\begin{scelta}{Come scegliere la formula}
\renewcommand{\arraystretch}{1.25}
\begin{tabular}{@{}P{0.56\linewidth}P{0.38\linewidth}@{}}
  \toprule
  \textbf{Se il testo\ldots} & \textbf{si usa\ldots} \\
  \midrule
  descrive due (o più) numeri prodotti insieme dallo stesso esperimento & la pmf congiunta \vv{congiunta} \\
  chiede la pmf di una sola variabile, partendo dalla congiunta & le marginali \vv{marginali} \\
  chiede se le variabili sono indipendenti & la fattorizzazione \vv{indipvar}, \vv{indippiu} \\
  chiede la distribuzione di $Y$ sapendo che $X = x$ & la condizionata \vv{condpmf} \\
  dà $p_X$ e il «canale» $p_{Y \mid X}$ & congiunta = prodotto \vv{condpmf}, poi $p_Y$ \vv{margcond} e, per
    risalire a $X$, Bayes \vv{bayespmf} \\
  costruisce tre o più variabili un passo alla volta & la regola della catena \vv{catena} \\
  \bottomrule
\end{tabular}

Il modo più ordinato di lavorare è scrivere la congiunta come tabella, con le marginali a margine: tutte le altre
grandezze si leggono da lì.
\end{scelta}

\scheda{congiunta}{
  \usoquando{L'esperimento produce più numeri insieme, spesso legati fra loro: il bit trasmesso e quello ricevuto, i
    risultati di due lanci, altezza e peso di una persona.}
  \usocome{Si scrive una tabella con una riga per ogni $x$ e una colonna per ogni $y$; si controlla che i numeri
    siano $\geq 0$ e sommino a 1.}
  \usoesempio{Due bit con $p_{X,Y}(0, 0) = \frac{1}{3}$, $p_{X,Y}(0, 1) = \frac{2}{9}$, $p_{X,Y}(1, 0) = \frac{1}{9}$,
    $p_{X,Y}(1, 1) = \frac{1}{3}$: sono positivi e sommano a $\frac{3 + 2 + 1 + 3}{9} = 1$.}
}

\scheda{marginali}{
  \usoquando{Si conosce la congiunta e serve la pmf di una sola delle variabili.}
  \usocome{Si somma sulla variabile che non interessa: i totali delle righe danno $p_X$, quelli delle colonne $p_Y$.}
  \usoattenzione{Dalle marginali, in generale, non si ricava la congiunta: due bit sempre uguali e due bit
    indipendenti hanno le stesse marginali ($\frac{1}{2}, \frac{1}{2}$) ma congiunte diverse.}
  \usoesempio{La tabella dell'esempio precedente: $p_X(0) = \frac{1}{3} + \frac{2}{9} = \frac{5}{9}$,
    $p_X(1) = \frac{4}{9}$; $p_Y(0) = \frac{1}{3} + \frac{1}{9} = \frac{4}{9}$, $p_Y(1) = \frac{5}{9}$.}
}

\scheda{indipvar}{
  \usoquando{Per \emph{verificare} se due variabili sono indipendenti, o per \emph{costruire} la congiunta di due
    variabili che lo sono (lanci ripetuti, bit indipendenti): in quel caso le marginali bastano.}
  \usocome{Si controlla che ogni casella sia il prodotto dei suoi due totali a margine.}
  \usoattenzione{Basta una sola casella che non si fattorizza perché le variabili siano dipendenti.}
  \usoesempio{Due lanci di un dado: $\frac{1}{36} = \frac{1}{6} \cdot \frac{1}{6}$ in ogni casella, indipendenti.
    I due bit di prima: $p_{X,Y}(0, 0) = \frac{1}{3} = \frac{27}{81}$, ma $p_X(0)\, p_Y(0) = \frac{5}{9} \cdot
    \frac{4}{9} = \frac{20}{81}$: dipendenti.}
}

\scheda{condpmf}{
  \usoquando{Si vuole la distribuzione di una variabile sapendo il valore dell'altra; oppure si conosce un
    «meccanismo» $p_{Y \mid X}$ (un canale, una regola) e serve la congiunta.}
  \usocome{Si divide la riga (o la colonna) della congiunta per il suo totale a margine. Al contrario, congiunta =
    condizionata $\times$ marginale.}
  \usoesempio{Canale binario: $p_X(0) = 0{,}6$, $p_X(1) = 0{,}4$, bit invertito con probabilità $0{,}1$. La congiunta
    è $p_{X,Y}(0, 0) = 0{,}9 \cdot 0{,}6 = 0{,}54$, $p_{X,Y}(0, 1) = 0{,}06$, $p_{X,Y}(1, 0) = 0{,}04$,
    $p_{X,Y}(1, 1) = 0{,}36$.}
}

\scheda{margcond}{
  \usoquando{Si conoscono la marginale di una variabile e la condizionata dell'altra, e serve la marginale dell'altra:
    è la probabilità totale scritta con le pmf.}
  \usoesempio{Canale binario: $p_Y(1) = 0{,}1 \cdot 0{,}6 + 0{,}9 \cdot 0{,}4 = 0{,}42$ e $p_Y(0) = 0{,}58$, cioè i
    totali delle colonne della congiunta.}
}

\scheda{bayespmf}{
  \usoquando{Si osserva $Y$ e si vuole risalire a $X$: il bit trasmesso dato quello ricevuto, la causa dato l'effetto.}
  \usocome{Numeratore: la casella $(x, y)$ della congiunta. Denominatore: il totale della colonna $y$.}
  \usoesempio{Canale binario: $p_{X \mid Y}(1 \mid 1) = \frac{0{,}36}{0{,}42} = \frac{6}{7}$ e
    $p_{X \mid Y}(0 \mid 1) = \frac{0{,}06}{0{,}42} = \frac{1}{7}$.}
}

\scheda{catena}{
  \usoquando{Si costruisce la congiunta di tre o più variabili un passo alla volta, come in un albero: prima una
    variabile, poi la seconda sapendo la prima, poi la terza sapendo le prime due.}
  \usocome{Si moltiplicano le probabilità lungo il percorso dell'albero; l'ordine delle variabili si sceglie in modo
    che ogni condizionata sia nota.}
  \usoesempio{$P(Z = 1) = 0{,}5$, $P(Y = 1 \mid Z = 1) = 0{,}8$, $P(X = 1 \mid Y = 1, Z = 1) = 0{,}9$:
    $p_{X,Y,Z}(1, 1, 1) = 0{,}5 \cdot 0{,}8 \cdot 0{,}9 = 0{,}36$.}
}

\scheda{indippiu}{
  \usoquando{Si deve decidere se tre o più variabili sono indipendenti, o si costruisce la congiunta di più prove
    indipendenti (più bit, più lanci).}
  \usocome{Si controlla che \emph{tutta} la congiunta sia il prodotto delle marginali; nella regola della catena,
    ogni condizionata deve coincidere con la sua marginale.}
  \usoattenzione{L'indipendenza a coppie non basta. Due bit indipendenti con probabilità $\frac{1}{2}$ e il loro bit di
    parità $B_3 = B_1 \oplus B_2$: ogni coppia è indipendente (tutte le caselle valgono $\frac{1}{4}$), ma
    $p(1, 1, 1) = 0 \neq \frac{1}{8}$.}
  \usoesempio{Tre bit indipendenti, ognuno uguale a 1 con probabilità $\alpha$: $p(0, 1, 1) = (1 - \alpha)\,
    \alpha^2$, un fattore $\alpha$ per ogni 1 e $1 - \alpha$ per ogni 0.}
}
